% answers_short.tex —— 1999 年·普通高等学校招生全国统一考试·广东卷 · 简短答案速查
% 由 tools/make_short_answers.py 自 answers.tex 的 \gksolhead/\gkanswer 提取，
% 并与 1999/1999_数学真题_广东_答案要点.md 逐题比对后生成，勿手改（改源 MD 后重跑）。

\begin{gknotice}
\textbf{使用说明}：本节只给最终答案便于快速核对，完整推导见解析版。答案为编写者独立推导，\textbf{非官方标准答案}。
\end{gknotice}

\begin{gknotice}
选择题答案中的 U 表示待核，未给出猜测答案，详情见解析版。
\end{gknotice}

\gkpart{选择题}
\begin{center}\begin{tabular}{ccccccccccccc}
\toprule
题号 & 1 & 2 & 3 & 4 & 5 & 6 & 7 & 8 & 9 & 10 & 11 & 12 \\
\midrule
答案 & C & A & A & C & B & B & B & A & C & U & B & C \\
\bottomrule
\end{tabular}\end{center}

\gkpart{填空题}
\begin{itemize}[leftmargin=2.6em]
  \item[\textbf{13.}] 条件答案 $e=\dfrac12$；原题按可见字面未决。
  \item[\textbf{14.}] 12。
  \item[\textbf{15.}] $[9,+\infty)$。
  \item[\textbf{16.}] $m\perp\alpha,\ n\perp\beta,\ \alpha\perp\beta\Rightarrow m\perp n$；或 $m\perp n,\ m\perp\alpha,\ n\perp\beta\Rightarrow\alpha\perp\beta$。
\end{itemize}

\gkpart{解答题}
\begin{itemize}[leftmargin=2.6em]
  \item[\textbf{17.}] $x=100$。
  \item[\textbf{18.}] 最大值 $\dfrac{\sqrt6}{12}$，对应 $\theta=\arctan\sqrt{\dfrac32}$。
  \item[\textbf{19.}] $\dfrac{a^2}{\sqrt2}$；$\sqrt2a$；$\dfrac{a^3}{2\sqrt2}$。
  \item[\textbf{20.}] 当 $0<\beta<\alpha$、$0<r_0<1$ 时，至少 $\left\lceil\dfrac{\lg(\beta/\alpha)}{\lg(1-r_0)}\right\rceil$ 对；$L_1=3125\ \mathrm{mm}$，$L_2=2500\ \mathrm{mm}$，$L_3=2000\ \mathrm{mm}$。
  \item[\textbf{21.}] 条件性结论：若原定义本意为 $f(x_n)=n$，则（Ⅰ）$x_n=\sum_{k=0}^{n-1}b^{-k}=\dfrac{1-b^{-n}}{1-b^{-1}}$；（Ⅱ）$b>1$ 时极限为 $\dfrac{b}{b-1}$，$0<b<1$ 时极限为 $+\infty$；（Ⅲ）在 $[x_n,x_{n+1}]$ 上 $f(x)=n+b^n(x-x_n)$，定义域分别为 $[0,\dfrac{b}{b-1})$ 与 $[0,+\infty)$。原页字面为 $x(x_n)=n$ 且小问（Ⅰ）列式异常，因此以上不是无条件答案。
  \item[\textbf{22.}] $(1-a)x^2-2ax+(a+1)y^2=0$，实际轨迹取 $0\leq x<a$；$0<a<1$ 时为椭圆，$a>1$ 时为双曲线的相应部分。
\end{itemize}
